\section*{Приложения}
\addcontentsline{toc}{section}{Приложения}
\label{sec:Appendix} \index{Appendix}

В приложениях приведены ключевые артефакты реализации, необходимые для воспроизведения экспериментов: системные промпты, схемы структурированного вывода, конфигурация обучения, примеры данных и реализация процедуры валидации графа.

\subsection*{Приложение А. Системный промпт RuleTaker (графовый режим)}
\addcontentsline{toc}{subsection}{Приложение А. Системный промпт RuleTaker}

Системный промпт задаёт жёсткий формат вывода \texttt{<think>...</think><answer>...</answer>}, в котором блок \texttt{<answer>} содержит JSON-массив рёбер графа рассуждения и финальную метку.

\begin{lstlisting}[basicstyle=\footnotesize\ttfamily]
You are a logical reasoning assistant. Given a context containing facts
and rules about entities, determine whether a given statement can be
logically derived from the provided information.

Analyze the context step by step in <think> tags, then give your answer
in <answer> tags.

The <answer> block MUST contain:
1. A JSON array of reasoning edges -- each edge is ["evidence",
   "conclusion"], meaning "evidence supports conclusion".
2. On a separate line: exactly 'entailment' or 'not entailment'.

Rules for building edges:
- Each edge is ["A", "B"] where A is a fact or derived claim that
  supports B.
- Facts (from the context) are leaves -- they appear only as evidence,
  never as conclusion.
- The statement you are checking is the root -- it appears only as
  conclusion, never as evidence.
- Write claims as "subject relation object" (e.g. "cow chases lion").
- For entailment: pick ONE rule per step and include ALL its conditions.
- For not entailment: show that NO rule works -- for each possible rule,
  show ONE condition that fails.
- Do NOT wrap the JSON in markdown code fences.

Example:
<think>cow chases lion is a fact. Rule says: if cow chases lion then
lion needs tiger. So lion needs tiger.</think>
<answer>
[["cow chases lion", "lion needs tiger"]]
entailment
</answer>
\end{lstlisting}

\subsection*{Приложение Б. Схемы структурированного вывода (Pydantic)}
\addcontentsline{toc}{subsection}{Приложение Б. Схемы структурированного вывода}

Схемы Pydantic задают контракт извлечения графа вспомогательной LLM (\texttt{StructuredGraphCreator}) и гарантируют синтаксически корректный JSON-выход.

\begin{lstlisting}[language=Python, basicstyle=\footnotesize\ttfamily]
class RuleTakerFact(BaseModel):
    src: str
    rel: str
    tgt: str
    negated: bool = False

class RuleTakerSimplePath(BaseModel):
    path: List[RuleTakerFact]

class ReasoningGraphNode(BaseModel):
    node_id: int
    type: str  # question | fact | condition | rule
    src: Optional[str] = None
    rel: Optional[str] = None
    tgt: Optional[str] = None
    fact_id: Optional[str] = None
    rule_id: Optional[str] = None
    condition_id: Optional[int] = None

class ReasoningGraphConnection(BaseModel):
    from_node: int
    to_node: int

class ReasoningGraph(BaseModel):
    nodes: ReasoningGraphNodes
    adj_lists: ReasoningGraphConnections
\end{lstlisting}

\subsection*{Приложение В. Конфигурация обучения GRPO (VERL)}
\addcontentsline{toc}{subsection}{Приложение В. Конфигурация обучения GRPO}

Фрагмент команды запуска основной серии экспериментов на RuleTaker (платформа VERL, одна GPU класса H100).

\begin{lstlisting}[basicstyle=\footnotesize\ttfamily]
python3 -m verl.trainer.main_ppo \
  algorithm.adv_estimator=grpo \
  algorithm.kl_ctrl.kl_coef=0.1 \
  data.train_batch_size=128 \
  data.max_prompt_length=1024 \
  data.max_response_length=4096 \
  actor_rollout_ref.model.path=Qwen/Qwen3-0.6B \
  actor_rollout_ref.model.lora_rank=64 \
  actor_rollout_ref.model.lora_alpha=16 \
  actor_rollout_ref.model.enable_gradient_checkpointing=True \
  actor_rollout_ref.actor.optim.lr=1e-4 \
  actor_rollout_ref.actor.ppo_mini_batch_size=64 \
  actor_rollout_ref.actor.use_kl_loss=True \
  actor_rollout_ref.actor.kl_loss_coef=0.001 \
  actor_rollout_ref.rollout.name=vllm \
  actor_rollout_ref.rollout.gpu_memory_utilization=0.7 \
  actor_rollout_ref.rollout.n=8 \
  actor_rollout_ref.rollout.temperature=0.7 \
  reward.custom_reward_function.path=verl_exp/edge_reasoning_reward.py \
  +reward.custom_reward_function.reward_kwargs.graph_validity_scale=1.0 \
  +reward.custom_reward_function.reward_kwargs.graph_validity_reward=100.0 \
  +reward.custom_reward_function.reward_kwargs.correct_answer_reward=100.0 \
  trainer.total_epochs=1 trainer.save_freq=10 trainer.test_freq=10
\end{lstlisting}

\subsection*{Приложение Г. Пример обучающего примера RuleTaker}
\addcontentsline{toc}{subsection}{Приложение Г. Пример данных}

Каждый пример датасета \texttt{KGReasoning} содержит промпт, эталонную метку и отображение \texttt{proof\_rules} (заключение $\mapsto$ списки допустимых наборов посылок).

\begin{lstlisting}[basicstyle=\footnotesize\ttfamily]
{
  "data_source": "ruletaker_kg",
  "ability": "logical_reasoning",
  "extra_info": {
    "question": "the lion needs the tiger",
    "context": "The cow chases the lion. If something chases the lion
                then the lion needs the tiger. ..."
  },
  "reward_model": {
    "style": "rule",
    "ground_truth": {
      "label": "entailment",
      "proof_rules": {
        "lion needs tiger": [["cow chases lion"]]
      }
    }
  }
}
\end{lstlisting}

\subsection*{Приложение Д. Процедура валидации подграфа}
\addcontentsline{toc}{subsection}{Приложение Д. Процедура валидации подграфа}

Ключевая процедура, формирующая обучающий сигнал: рекурсивный обход узлов-заключений от вопроса с проверкой полноты набора посылок (раздел~\ref{subsec:metrics-formalization}).

\begin{lstlisting}[language=Python, basicstyle=\footnotesize\ttfamily]
def check_node(node):
    if node in visited:
        return True
    visited.add(node); total_nodes += 1
    evidences = by_conclusion.get(node)
    if evidences is None:          # leaf fact -- valid
        valid_nodes += 1; return True
    gt_variants = norm_rules.get(node)
    if gt_variants is None:        # no ground-truth rule for node
        return False
    if any(evidences == variant for variant in gt_variants):
        valid_nodes += 1           # fully valid (score 1.0)
        for ev in evidences: check_node(ev)
        return True
    valid_nodes += 0.5             # partially valid (score 0.5)
    for ev in evidences: check_node(ev)
    return False

check_node(norm_question)
coverage = valid_nodes / total_nodes if total_nodes > 0 else 0.0
\end{lstlisting}

\subsection*{Приложение Е. Baseline-промпт и пользовательские шаблоны}
\addcontentsline{toc}{subsection}{Приложение Е. Baseline-промпт и шаблоны}

Для контрастного baseline-режима ($\alpha=0$) используется упрощённый системный промпт, не требующий построения графа рёбер, а только финальной метки.

\begin{lstlisting}[basicstyle=\footnotesize\ttfamily]
# SYSTEM_PROMPT_BASELINE
You are a logical reasoning assistant. Given a context containing facts
and rules about entities, determine whether a given statement can be
logically derived from the provided information.

Analyze the context step by step, identifying which facts and rules are
relevant to prove or disprove the statement. Put your reasoning in
<think> tags and your final answer in <answer> tags. Your answer must be
exactly "entailment" if the statement follows from the context, or
"not entailment" if it does not.

# USER_TEMPLATE (graph mode)
Context:
{context}

Statement: {question}

Can this statement be derived from the facts and rules in the context?
Reply with <think>...</think><answer>[[edges...]]\n
entailment OR not entailment</answer>.
\end{lstlisting}

\subsection*{Приложение Ж. Вычисление меры сходства графов на основе GED}
\addcontentsline{toc}{subsection}{Приложение Ж. Вычисление GES}

Нормированная мера Graph Edit Similarity вычисляется на основе расстояния редактирования графов (\texttt{networkx}) с единичными стоимостями операций над вершинами и рёбрами.

\begin{lstlisting}[language=Python, basicstyle=\footnotesize\ttfamily]
def graph_edit_similarity(pred, gt):
    if pred.number_of_nodes() == 0 and gt.number_of_nodes() == 0:
        return 1.0
    ged = nx.graph_edit_distance(
        pred, gt,
        node_match=lambda a, b: a["label"] == b["label"],
        edge_match=lambda a, b: a["label"] == b["label"],
    )
    denom = max(pred.number_of_nodes() + pred.number_of_edges(),
                gt.number_of_nodes() + gt.number_of_edges())
    return 1.0 - ged / denom if denom > 0 else 0.0
\end{lstlisting}

\subsection*{Приложение З. Композиция функции награды}
\addcontentsline{toc}{subsection}{Приложение З. Композиция функции награды}

Итоговая награда складывается из шести компонент (формула~(\ref{eq:reward-total})). Ниже приведён фрагмент вычисления, отражающий обработку формата, корректности метки и графового покрытия, включая частный случай отрицательного примера.

\begin{lstlisting}[language=Python, basicstyle=\footnotesize\ttfamily]
# tags: number of <think> / <answer> blocks
reward_think_tags = think_tag_reward if len(tags["think"]) == 1 else \
    (-(len(tags["think"]) - 1) * think_tag_reward if tags["think"] else 0.0)
reward_answer_tags = answer_tag_reward if len(tags["answer"]) == 1 else \
    (-answer_tag_reward * 10 if not tags["answer"]
     else -(len(tags["answer"]) - 1) * answer_tag_reward)

# parse the JSON edge list out of <answer>
predicted_edges, label_text, parse_meta = parse_answer_block(answer_text)
reward_graph_parse = graph_parse_bonus if parse_meta["graph_parse_ok"] else 0.0

# final-label correctness
normalized = normalize_answer(label_text)
is_correct = normalized is not None and normalized == gt_label
reward_correct_answer = correct_answer_reward if is_correct else 0.0

# graph coverage (dense structural signal)
if gt_label == "not entailment":
    graph_coverage = 1.0 if not predicted_edges else 0.5
elif gt_proof_rules:
    graph_coverage, _ = validate_subgraph(
        predicted_edges, gt_proof_rules, question_text)
else:
    graph_coverage = 1.0 if not predicted_edges else 0.0

reward_graph_validity = (
    graph_coverage * graph_validity_reward * graph_validity_scale)  # alpha
reward_format = -min(content_len * format_penalty, 100.0)

reward = (reward_think_tags + reward_answer_tags + reward_graph_parse
          + reward_correct_answer + reward_graph_validity + reward_format)
\end{lstlisting}
